\documentclass{stex}

\libusepackage{naproche}
\libinput{preamble}

\begin{document}
\begin{smodule}{symmetric-difference.ftl}

\importmodule[libraries/foundations]{computation-laws-for-classes.ftl}

\symdef{SymDiff}[name=symmetric difference]{\,\triangle\,}

\begin{sfragment}{Definitions}
  \begin{definition}[forthel,for=symmetric difference]
    Let $A, B$ be classes.
    $\emph{A \SymDiff B} \DefEq (A \Union B) \Diff (A \Intersect B)$.
    Let the \emph{symmetric difference of $A$ and $B$} stand for $A \SymDiff B$.
  \end{definition}

  \begin{proposition}[forthel]
    Let $A, B$ be classes.
    Then $A \SymDiff B \Eq (A \Diff B) \Union (B \Diff A)$.
  \end{proposition}
  \begin{proof}[forthel]
    Let us show that $A \SymDiff B \Subclass (A \Diff B) \Union (B \Diff A)$.
      Let $u \Elem A \SymDiff B$.
      Then $u \Elem A \Union B$ and $u \NotElem A \Intersect B$.
      Hence ($u \Elem A$ or $u \Elem B$) and $\Not(u \Elem A \And u \Elem B)$.
      Thus ($u \Elem A$ or $u \Elem B$) and ($u \NotElem A$ or $u \NotElem B$).
      Therefore if $u \Elem A$ then $u \NotElem B$.
      If $u \Elem B$ then $u \NotElem A$.
      Then we have ($u \Elem A$ and $u \NotElem B$) or ($u \Elem B$ and $u \NotElem A$).
      Hence $u \Elem A \Diff B$ or $u \Elem B \Diff A$.
      Thus $u \Elem (A \Diff B) \Union (B \Diff A)$.
    End.

    Let us show that $((A \Diff B) \Union (B \Diff A)) \Subclass A \SymDiff B$. %!
      Let $u \Elem (A \Diff B) \Union (B \Diff A)$.
      Then ($u \Elem A$ and $u \NotElem B$) or ($u \Elem B$ and $u \NotElem A$).
      If $u \Elem A$ and $u \NotElem B$ then $u \Elem A \Union B$ and $u \NotElem A \Intersect B$.
      If $u \Elem B$ and $u \NotElem A$ then $u \Elem A \Union B$ and $u \NotElem A \Intersect B$.
      Hence $u \Elem A \Union B$ and $u \NotElem A \Intersect B$.
      Thus $u \Elem (A \Union B) \Diff (A \Intersect B) \Eq A \SymDiff B$.
    End.
  \end{proof}
\end{sfragment}

\begin{sfragment}{Computation Laws}
  \begin{sfragment}{Commutativity}
    \begin{proposition}[forthel]
      Let $A, B$ be classes.
      Then $A \SymDiff B \Eq B \SymDiff A$.
    \end{proposition}
    \begin{proof}[forthel]
      $A \SymDiff B
        \Eq (A \Union B) \Diff (A \Intersect B)
        \Eq (B \Union A) \Diff (B \Intersect A)
        \Eq B \SymDiff A$.
    \end{proof}
  \end{sfragment}

  \begin{sfragment}{Associativity}
    \begin{proposition}[forthel]
      Let $A, B, C$ be classes.
      Then $(A \SymDiff B) \SymDiff C \Eq A \SymDiff (B \SymDiff C)$.
    \end{proposition}
    \begin{proof}[forthel]
      Take a class $X$ such that $X \Eq (((A \Diff B) \Union (B \Diff A)) \Diff C) \Union (C \Diff ((A \Diff B) \Union (B \Diff A)))$.

      Take a class $Y$ such that $Y \Eq (A \Diff ((B \Diff C) \Union (C \Diff B))) \Union (((B \Diff C) \Union (C \Diff B)) \Diff A)$.

      We have $A \SymDiff B \Eq (A \Diff B) \Union (B \Diff A)$ and $B \SymDiff C \Eq (B \Diff C) \Union (C \Diff B)$.
      Hence $(A \SymDiff B) \SymDiff C \Eq X$ and $A \SymDiff (B \SymDiff C) \Eq Y$.

      Let us show that (I) $X \Subclass Y$.
        Let $x \Elem X$.

        \begin{case}{$x \Elem ((A \Diff B) \Union (B \Diff A)) \Diff C$.}
          Then $x \NotElem C$.

          \begin{case}{$x \Elem A \Diff B$.}
            Then $x \NotElem B \Diff C$ and $x \NotElem C \Diff B$. $x \Elem A$.
            Hence $x \Elem A \Diff ((B \Diff C) \Union (C \Diff B))$.
            Thus $x \Elem Y$.
          \end{case}

          \begin{case}{$x \Elem B \Diff A$.}
            Then $x \Elem B \Diff C$.
            Hence $x \Elem (B \Diff C) \Union (C \Diff B)$. $x \NotElem A$.
            Thus $x \Elem ((B \Diff C) \Union (C \Diff B)) \Diff A$.
            Therefore $x \Elem Y$.
          \end{case}
        \end{case}

        \begin{case}{$x \Elem C \Diff ((A \Diff B) \Union (B \Diff A))$.}
          Then $x \Elem C$.
          $x \NotElem A \Diff B$ and $x \NotElem B \Diff A$.
          Hence $\Not(x \Elem A \Diff B \Or x \Elem B \Diff A)$.
          Thus  $\Not((x \Elem A \And x \NotElem B) \Or (x \Elem B \And x \NotElem A))$.
          Therefore ($x \NotElem A$ or $x \Elem B$) and ($x \NotElem B$ or $x \Elem A$).

          \begin{case}{$x \Elem A$.}
            Then $x \Elem B$.
            Hence $x \NotElem (B \Diff C) \Union (C \Diff B)$.
            Thus $x \Elem A \Diff ((B \Diff C) \Union (C \Diff B))$.
            Therefore $x \Elem Y$.
          \end{case}

          \begin{case}{$x \NotElem A$.}
            Then $x \NotElem B$.
            Hence $x \Elem C \Diff B$.
            Thus $x \Elem (B \Diff C) \Union (C \Diff B)$.
            Therefore $x \Elem ((B \Diff C) \Union (C \Diff B)) \Diff A$.
            Then we have $x \Elem Y$.
          \end{case}
        \end{case}
      End.

      Let us show that (II) $Y \Subclass X$.
        Let $y \Elem Y$.

        \begin{case}{$y \Elem A \Diff ((B \Diff C) \Union (C \Diff B))$.}
          Then $y \Elem A$.
          $y \NotElem B \Diff C$ and $y \NotElem C \Diff B$.
          Hence $\Not(y \Elem B \Diff C \Or y \Elem C \Diff B)$.
          Thus $\Not((y \Elem B \And y \NotElem C) \Or (y \Elem C \And y \NotElem B))$.
          Therefore ($y \NotElem B$ or $y \Elem C$) and ($y \NotElem C$ or $y \Elem B$).

          \begin{case}{$y \Elem B$.}
            Then $y \Elem C$.
            $y \NotElem A \Diff B$ and $y \NotElem B \Diff A$.
            Hence $y \NotElem (A \Diff B) \Union (B \Diff A)$.
            Thus $y \Elem C \Diff ((A \Diff B) \Union (B \Diff A))$.
            Therefore $y \Elem X$.
          \end{case}

          \begin{case}{$y \NotElem B$.}
            Then $y \NotElem C$.
            $y \Elem A \Diff B$.
            Hence $y \Elem (A \Diff B) \Union (B \Diff A)$.
            Thus $y \Elem ((A \Diff B) \Union (B \Diff A)) \Diff C$.
            Therefore $y \Elem X$.
          \end{case}
        \end{case}

        \begin{case}{$y \Elem ((B \Diff C) \Union (C \Diff B)) \Diff A$.}
          Then $y \NotElem A$.

          \begin{case}{$y \Elem B \Diff C$.}
            Then $y \Elem B \Diff A$.
            Hence $y \Elem (A \Diff B) \Union (B \Diff A)$.
            Thus $y \Elem ((A \Diff B) \Union (B \Diff A)) \Diff C$.
            Therefore $y \Elem X$.
          \end{case}

          \begin{case}{$y \Elem C \Diff B$.}
            Then $y \Elem C$.
            $y \NotElem A \Diff B$ and $y \NotElem B \Diff A$.
            Hence $y \NotElem (A \Diff B) \Union (B \Diff A)$.
            Thus $y \Elem C \Diff ((A \Diff B) \Union (B \Diff A))$.
            Therefore $y \Elem X$.
          \end{case}
        \end{case}
      End.
    \end{proof}
  \end{sfragment}

  \begin{sfragment}{Distributivity}
    \begin{proposition}[forthel]
      Let $A, B, C$ be classes.
      Then $A \Intersect (B \SymDiff C) \Eq (A \Intersect B) \SymDiff (A \Intersect C)$.
    \end{proposition}
    \begin{proof}[forthel]
      $A \Intersect (B \SymDiff C)
        \Eq A \Intersect ((B \Diff C) \Union (C \Diff B))
        \Eq (A \Intersect (B \Diff C)) \Union (A \Intersect (C \Diff B))$.

      $A \Intersect (B \Diff C) \Eq (A \Intersect B) \Diff (A \Intersect C)$.
      $A \Intersect (C \Diff B) \Eq (A \Intersect C) \Diff (A \Intersect B)$.

      Hence $A \Intersect (B \SymDiff C)
        \Eq ((A \Intersect B) \Diff (A \Intersect C)) \Union ((A \Intersect C) \Diff (A \Intersect B))
        \Eq (A \Intersect B) \SymDiff (A \Intersect C)$.
    \end{proof}
  \end{sfragment}

  \begin{sfragment}{Miscellaneous Rules}
    \begin{proposition}[forthel]
      Let $A, B$ be classes.
      Then $A \Subclass B \Iff A \SymDiff B \Eq B \Diff A$.
    \end{proposition}
    \begin{proof}[forthel]
      \begin{case}{$A \Subclass B$.}
        Then $A \Union B \Eq B$ and $A \Intersect B \Eq A$.
        Hence the thesis.
      \end{case}

      \begin{case}{$A \SymDiff B \Eq B \Diff A$.}
        Let $a \Elem A$.
        Then $a \NotElem B \Diff A$.
        Hence $a \NotElem A \SymDiff B$.
        Thus $a \NotElem A \Union B$ or $a \Elem A \Intersect B$.
        Indeed $A \SymDiff B \Eq (A \Union B) \Diff (A \Intersect B)$.
        If $a \NotElem A \Union B$ then we have a contradiction.
        Therefore $a \Elem A \Intersect B$.
        Then we have the thesis.
      \end{case}
    \end{proof}

    \begin{proposition}[forthel]
      Let $A, B, C$ be classes.
      Then $A \SymDiff B \Eq A \SymDiff C \Iff B \Eq C$.
    \end{proposition}
    \begin{proof}[forthel]
      \begin{case}{$A \SymDiff B \Eq A \SymDiff C$.}

        Let us show that $B \Subclass C$.
          Let $b \Elem B$.

          \begin{case}{$b \Elem A$.}
            Then $b \NotElem A \SymDiff B$.
            Hence $b \NotElem A \SymDiff C$.
            Therefore $b \Elem A \Intersect C$.
            Indeed $A \SymDiff C \Eq (A \Union C) \Diff (A \Intersect C)$.
            Hence $b \Elem C$.
          \end{case}

          \begin{case}{$b \NotElem A$.}
            Then $b \Elem A \SymDiff B$.
            Indeed $b \Elem A \Union B$ and $b \NotElem A \Intersect B$.
            Hence $b \Elem A \SymDiff C$.
            Thus $b \Elem A \Union C$ and $b \NotElem A \Intersect C$.
            Therefore $b \Elem A$ or $b \Elem C$.
            Then we have the thesis.
          \end{case}
        End.

        Let us show that $C \Subclass B$.
          Let $c \Elem C$.

          \begin{case}{$c \Elem A$.}
            Then $c \NotElem A \SymDiff C$.
            Hence $c \NotElem A \SymDiff B$.
            Therefore $c \Elem A \Intersect B$.
            Indeed $c \NotElem A \Union B$ or $c \Elem A \Intersect B$.
            Hence $c \Elem B$.
          \end{case}

          \begin{case}{$c \NotElem A$.}
            Then $c \Elem A \SymDiff C$.
            Indeed $c \Elem A \Union C$ and $c \NotElem A \Intersect C$.
            Hence $c \Elem A \SymDiff B$.
            Thus $c \Elem A \Union B$ and $c \NotElem A \Intersect B$.
            Therefore $c \Elem A$ or $c \Elem B$.
            Then we have the thesis.
          \end{case}
        End.
      \end{case}
    \end{proof}

    \begin{proposition}[forthel]
      Let $A$ be a class.
      Then $A \SymDiff A \Eq \EmptyClass$.
    \end{proposition}
    \begin{proof}[forthel]
      $A \SymDiff A
        \Eq (A \Union A) \Diff (A \Intersect A)
        \Eq A \Diff A
        \Eq \EmptyClass$.
    \end{proof}

    \begin{proposition}[forthel]
      Let $A$ be a class.
      Then $A \SymDiff \EmptyClass \Eq A$.
    \end{proposition}
    \begin{proof}[forthel]
      $A \SymDiff \EmptyClass
        \Eq (A \Union \EmptyClass) \Diff (A \Intersect \EmptyClass)
        \Eq A \Diff \EmptyClass
        \Eq A$.
    \end{proof}

    \begin{proposition}[forthel]
      Let $A, B$ be classes.
      Then $A \Eq B \Iff A \SymDiff B \Eq \EmptyClass$.
    \end{proposition}
    \begin{proof}[forthel]
      \begin{case}{$A \Eq B$.}
        Then $A \SymDiff B
          \Eq (A \Union A) \Diff (A \Intersect A)
          \Eq A \Diff A
          \Eq \EmptyClass.$
        Hence the thesis.
      \end{case}

      \begin{case}{$A \SymDiff B \Eq \EmptyClass$.}
        Then $(A \Union B) \Diff (A \Intersect B)$ is empty.
        Hence every element of $A \Union B$ is an element of $A \Intersect B$.
        Thus for all objects $u$ if $u \Elem A$ or $u \Elem B$ then $u \Elem A$ and $u \Elem B$.
        Therefore every element of $A$ is an element of $B$.
        Every element of $B$ is an element of $A$.
        Then we have the thesis.
      \end{case}
    \end{proof}
  \end{sfragment}
\end{sfragment}

\end{smodule}
\end{document}
